\documentclass[10pt,a4paper]{amsart}
\usepackage{amsmath,amssymb,amsthm,graphicx}
% Retained preamble for the Stacks Project chapter snapshot in this directory.
% The chapter's own preamble.tex is a copy of the shared Stacks Project preamble; the block
% below is sliced MECHANICALLY, line for line, from it. Nothing was added, removed or
% redefined; the wrapper class, the xr-hyper/hyperref wiring and the externaldocument list
% are omitted because the assembler supplies the document class.
%
% lines taken from preamble.tex: 3, 5, 9, 11, 12, 156, 157, 158, ... (51 lines total)
%
% package coverage restored from the chapter's own preamble (2026-10-13): xy, xy-2cell

\usepackage[all]{xy}
\xyoption{2cell}
\UseAllTwocells

\newenvironment{reference}{\comment}{\endcomment}
\newenvironment{slogan}{\comment}{\endcomment}
\newenvironment{history}{\comment}{\endcomment}
\theoremstyle{plain}
\newtheorem{theorem}[subsection]{Theorem}
\newtheorem{proposition}[subsection]{Proposition}
\newtheorem{lemma}[subsection]{Lemma}
\theoremstyle{definition}
\newtheorem{definition}[subsection]{Definition}
\newtheorem{example}[subsection]{Example}
\newtheorem{exercise}[subsection]{Exercise}
\newtheorem{situation}[subsection]{Situation}
\theoremstyle{remark}
\newtheorem{remark}[subsection]{Remark}
\newtheorem{remarks}[subsection]{Remarks}
\numberwithin{equation}{subsection}
\def\lim{\mathop{\mathrm{lim}}\nolimits}
\def\colim{\mathop{\mathrm{colim}}\nolimits}
\def\Spec{\mathop{\mathrm{Spec}}}
\def\Hom{\mathop{\mathrm{Hom}}\nolimits}
\def\Ext{\mathop{\mathrm{Ext}}\nolimits}
\def\SheafHom{\mathop{\mathcal{H}\!\mathit{om}}\nolimits}
\def\SheafExt{\mathop{\mathcal{E}\!\mathit{xt}}\nolimits}
\def\Sch{\mathit{Sch}}
\def\Mor{\mathop{\mathrm{Mor}}\nolimits}
\def\Ob{\mathop{\mathrm{Ob}}\nolimits}
\def\Sh{\mathop{\mathit{Sh}}\nolimits}
\def\NL{\mathop{N\!L}\nolimits}
\def\CH{\mathop{\mathrm{CH}}\nolimits}
\def\proetale{{pro\text{-}\acute{e}tale}}
\def\etale{{\acute{e}tale}}
\def\QCoh{\mathit{QCoh}}
\def\Ker{\mathop{\mathrm{Ker}}}
\def\Im{\mathop{\mathrm{Im}}}
\def\Coker{\mathop{\mathrm{Coker}}}
\def\Coim{\mathop{\mathrm{Coim}}}
\DeclareMathSymbol{\boxtimes}{\mathbin}{AMSa}{"02}
\def\QCohstack{\mathcal{QC}\!\mathit{oh}}
\def\Cohstack{\mathcal{C}\!\mathit{oh}}
\def\Spacesstack{\mathcal{S}\!\mathit{paces}}
\def\Quotfunctor{\mathrm{Quot}}
\def\Hilbfunctor{\mathrm{Hilb}}
\def\Curvesstack{\mathcal{C}\!\mathit{urves}}
\def\Polarizedstack{\mathcal{P}\!\mathit{olarized}}
\def\Complexesstack{\mathcal{C}\!\mathit{omplexes}}
\def\Pic{\mathop{\mathrm{Pic}}\nolimits}
\def\Picardstack{\mathcal{P}\!\mathit{ic}}
\def\Picardfunctor{\mathrm{Pic}}
\def\Deformationcategory{\mathcal{D}\!\mathit{ef}}
\usepackage{cleveref}
\crefname{originallemma}{original lemma}{original lemmas}
\Crefname{originallemma}{Original lemma}{Original lemmas}
\makeatletter
\newlabel{equation-beauville-laszlo-glueing-spaces}{{0F9N}{source}{}{sourcepointer-1}{}}
\newlabel{equation-beauville-laszlo-glueing-spaces@cref}{{[originallemma][0][]0F9N}{[1][0][]source}{}{}{}}
\makeatother
\crefname{originallemma}{original lemma}{original lemmas}
\Crefname{originallemma}{Original lemma}{Original lemmas}
\makeatletter
\newlabel{more-morphisms-lemma-quasi-finite-separated-pass-through-finite}{{05K0}{source}{}{sourcepointer-2}{}}
\newlabel{more-morphisms-lemma-quasi-finite-separated-pass-through-finite@cref}{{[originallemma][0][]05K0}{[1][0][]source}{}{}{}}
\makeatother
\crefname{originallemma}{original lemma}{original lemmas}
\Crefname{originallemma}{Original lemma}{Original lemmas}
\makeatletter
\newlabel{spaces-morphisms-lemma-properties-diagonal}{{03HK}{source}{}{sourcepointer-3}{}}
\newlabel{spaces-morphisms-lemma-properties-diagonal@cref}{{[originallemma][0][]03HK}{[1][0][]source}{}{}{}}
\makeatother
\title{The glueing functor is faithful on glueable spaces}
\author{The Stacks Project Authors}
\date{stacks-project:spaces-pushouts.tex:local-snapshot-20261012; GNU FDL 1.2 or later}
\begin{document}
\maketitle
\noindent\textit{Editorial note.} Statement span is lines 2647-2653 of sources/stacks-spaces-pushouts/pinned-original-source.tex; proof span is lines 2655-2683, the complete author proof (27 body lines).\par\medskip
\noindent\textit{Editorial note.} Navigation stubs: 3 cited-but-not-included labels are supplied as non-load-bearing pointers with their pinned tags where the snapshot tag index has them.\par\medskip
\noindent\textit{Editorial note.} extra\_packages is empty on purpose: the retained preamble already loads xy with options and a second declaration clashes in the assembler wrapper.\par\medskip
\noindent\textit{Editorial note.} Text is retained verbatim, including the author's own typos and placeholders.\par\medskip
\noindent\textit{Editorial note.} Queue index 15 of queue\_spaces-pushouts.json (pb=1725, stmt 2647-2653, proof 2655-2683).\par\medskip
\noindent\textit{Editorial note.} spaces-pushouts wave-1: fresh adjudication accept (2026-10-14); six-check dispatched pool (§2 row appended 02:4x).\par\medskip
\noindent\textit{Allowed standard prerequisites (editorial).}\begin{itemize}
\item Explicit prerequisite boundary. The cited inputs are prerequisites; the argument in the retained span is not replaced by any cited result.
\end{itemize}\medskip
\ifdefined\hypertarget\hypertarget{sourcepointer-1}{}\fi
\noindent\textit{Source reference.} Original lemma 0F9N: sources/stacks-spaces-pushouts/pinned-original-source.tex, source lines 2470--2474. Cited by the retained statement or proof but not part of the collected spans; supplied only as a navigation pointer. The argument does not depend on its content.\par\medskip
\ifdefined\hypertarget\hypertarget{sourcepointer-2}{}\fi
\noindent\textit{Source reference.} Original lemma 05K0: sources/stacks-spaces-pushouts/pinned-original-source.tex, source lines 2683--2683. Cited by the retained statement or proof but not part of the collected spans; supplied only as a navigation pointer. The argument does not depend on its content.\par\medskip
\ifdefined\hypertarget\hypertarget{sourcepointer-3}{}\fi
\noindent\textit{Source reference.} Original lemma 03HK: sources/stacks-spaces-pushouts/pinned-original-source.tex, source lines 2683--2683. Cited by the retained statement or proof but not part of the collected spans; supplied only as a navigation pointer. The argument does not depend on its content.\par\medskip
\begin{lemma}
\label{lemma-glueing-f}
Let $(R \to R', f)$ be a glueing pair, see above.
The functor (\ref{equation-beauville-laszlo-glueing-spaces})
is faithful on the full subcategory of
algebraic spaces $Y/X$ glueable for $(R \to R', f)$.
\end{lemma}

\begin{proof}
Let $f, g : Y \to Z$ be two morphisms of algebraic spaces over $X$
with $Y$ and $Z$ glueable for $(R \to R', f)$ such that $f$ and $g$ are mapped
to the same morphism in the category $\textit{Spaces}(U \leftarrow U' \to X')$.
We have to show the equalizer $E \to Y$ of $f$ and $g$ is an isomorphism.
Working \'etale locally on $Y$ we may assume $Y$ is an affine scheme.
Then $E$ is a scheme and the morphism $E \to Y$ is a monomorphism
and locally quasi-finite, see Morphisms of Spaces, Lemma
\ref{spaces-morphisms-lemma-properties-diagonal}.
Moreover, the base change of $E \to Y$ to $U$ and to $X'$ is
an isomorphism. As $Y$ is the disjoint union of the affine open
$V = U \times_X Y$ and the affine closed $V(f) \times_X Y$, we conclude
$E$ is the disjoint union of their isomorphic inverse images.
It follows in particular that $E$ is quasi-compact.
By Zariski's main theorem (More on Morphisms, Lemma
\ref{more-morphisms-lemma-quasi-finite-separated-pass-through-finite})
we conclude that $E$ is quasi-affine.
Set $B = \Gamma(E, \mathcal{O}_E)$ and $A = \Gamma(Y, \mathcal{O}_Y)$
so that we have an $R$-algebra homomorphism $A \to B$.
Since $E \to Y$ becomes an isomorphism after base change to $U$ and $X'$
we obtain ring maps $B \to A_f$ and $B \to A \otimes_R R'$
agreeing as maps into $A \otimes_R R'_f$. Since $A$ is glueable
for $(R \to R', f)$ we get a ring map $B \to A$ which is left inverse
to the map $A \to B$. The corresponding morphism $Y = \Spec(A) \to \Spec(B)$
maps into the open subscheme $E \subset \Spec(B)$ pointwise because
this is true after base change to $U$ and $X'$. Hence we get a morphism
$Y \to E$ over $Y$. Since $E \to Y$ is a monomorhism we conclude
$Y \to E$ is an isomorphism as desired.
\end{proof}

\end{document}
